\documentclass[11pt,aspectratio=169]{beamer}

\usetheme{default}
\setbeamertemplate{navigation symbols}{}
\setbeamertemplate{headline}{}
\setbeamertemplate{footline}[frame number]
\setbeamerfont{frametitle}{size=\large,series=\bfseries}

\usepackage{booktabs}
\usepackage{array}
\usepackage{xcolor}

\definecolor{qblue}{RGB}{30,80,160}
\definecolor{qlight}{RGB}{235,242,252}
\setbeamercolor{block title}{bg=qblue,fg=white}
\setbeamercolor{block body}{bg=qlight,fg=black}

\title{BDM Belief Incentive Compatibility}
\subtitle{Design review}
\author{Christina Sun}
\date{April 2026}

\begin{document}

\begin{frame}[plain]
\maketitle
\end{frame}

\begin{frame}{Literature}
Danz, Vesterlund, and Wilson (JEP 2024 and forthcoming WP) find that BDM fails behavioral incentive compatibility: 69\% of subjects pick $q = 0$ under p-BDM at induced $\theta = 0.2$. Their p-BDM methodology is not public.

\vspace{0.6em}
Chakraborty and Kendall (2025) formalize UJS. BDM admits many justifiable non-truthful reports; MPL admits one, may be better BIC.

\end{frame}

\begin{frame}{Research questions and contribution}
Research questions:
\begin{enumerate}
\item Why does BDM belief elicitation fail behavioral IC?
\item Does MPL perform better, despite theoretical equivalence?
\item Does the presentation format interact with task complexity for belief elicitation?
\end{enumerate}

\vspace{0.8em}
Contribution: identify the mechanism behind BDM's failure, test UJS's design implication empirically in belief elicitation, and measure the complexity-by-mechanism interaction. 
\end{frame}

\begin{frame}{Hypotheses}
\begin{description}
\item[H1.] BDM fails behavioral IC, via two weak conditions (DVW 2022):
\begin{itemize}
\item[H1a.] Explaining BDM's incentives does not increase accuracy.
\item[H1b.] Subjects do not choose the payoff-maximizer under pure incentives.
\end{itemize}
\item[H2.] Belief MPL achieves better behavioral IC than single-report BDM, despite theoretical equivalence, because MPL admits a unique justifiable action per row (UJS) while BDM admits many.
\item[H3.] The BDM-MPL truthful reporting gap increases with task complexity.
\end{description}
\end{frame}

\begin{frame}{Experimental design}
Ball and urn paradigm.
\vspace{0.7em}
\footnotesize
\setlength{\tabcolsep}{4pt}
\begin{tabular}{@{}p{0.03\linewidth} p{0.22\linewidth} p{0.32\linewidth} p{0.33\linewidth}@{}}
\toprule
T\# & Mechanism & Instructions & Identifies \\
\midrule
1 & Single-report BDM & Full info + quiz & H1 baseline \\
2 & Belief MPL & Full info + quiz & H2, format matters \\
3 & Single-report BDM & No info & H1a: info/no-info \\
4 & Pure-incentives test & --- & H1b: direct test \\
5 & Flat fee & Accuracy encouragement & No-mechanism benchmark \\
\bottomrule
\end{tabular}

\vspace{0.7em}
\small
Within-subject: priors (easy) and posteriors after signals (hard). Identifies H3.

Outcomes: Accuracy at $\varepsilon = 0$ and $\varepsilon = 5$pp (DVW 2022). Also multi-switching, response time, comprehension.

\end{frame}

\begin{frame}{Question 1: the overall design}
\begin{block}{}
Is there anything I am missing before we move forward?
\end{block}

\vspace{1em}
\begin{itemize}
\item Is the current design sufficient to identify H1 through H3?
\item How to design p-BDM pure incentive test? (coming up next)
\item Is the within-subject easy/hard manipulation enough for H3?
\end{itemize}
\end{frame}

\begin{frame}{p-BDM pure-incentives test (H1b)}
DVW report one p-BDM pure incentive test result: 69\% of subjects pick $q = 0$ at $\theta = 0.2$. WP is not available.

\vspace{0.7em}
\begin{itemize}
\item $q = 0$ pays with probability $0.5$ regardless of event.
\item $q = \theta = 0.2$ (the maximizer) pays with probability $0.68$ if $E$, $0.48$ otherwise, for an expected $0.52$.
\item The EV gap between the maximizer and the event-independent option is 2 percentage points.
\end{itemize}


\end{frame}

\begin{frame}{Proposal A: Example at $\theta = 0.2$}
\small
$P(\text{win} \mid E, q) = q + (1 - q^2)/2$, \quad $P(\text{win} \mid \lnot E, q) = (1 - q^2)/2$.

\vspace{0.5em}
\footnotesize
\begin{tabular}{@{}c c c c l@{}}
\toprule
Report $q$ & $P(\text{win} \mid E)$ & $P(\text{win} \mid \lnot E)$ & EV & \\
\midrule
0.0 & 0.500 & 0.500 & 0.500 & event-independent \\
0.1 & 0.595 & 0.495 & 0.515 & \\
0.2 & 0.680 & 0.480 & 0.520 & maximizer, $q^* = \theta$ \\
0.3 & 0.755 & 0.455 & 0.515 & \\
0.4 & 0.820 & 0.420 & 0.500 & \\
0.5 & 0.875 & 0.375 & 0.475 & \\
0.6 & 0.920 & 0.320 & 0.440 & \\
0.7 & 0.955 & 0.255 & 0.395 & \\
0.8 & 0.980 & 0.180 & 0.340 & \\
0.9 & 0.995 & 0.095 & 0.275 & \\
1.0 & 1.000 & 0.000 & 0.200 & \\
\bottomrule
\end{tabular}

\vspace{0.3em}
\footnotesize
Mechanism and $q$ values hidden. Subject compares event-contingent payoff pairs directly --- no contingent reasoning across $r$ is required.
\end{frame}

\begin{frame}{Proposal B: Example at $\theta = 0.2$}
Urn: 20 red balls, 80 blue. Event $E$ = red drawn. $P(E) = 0.2 = \theta$.

\vspace{0.6em}
Subject sees:

\vspace{0.3em}
\begin{quote}
You will report a number $q$ from $\{0.0, 0.1, 0.2, \ldots, 1.0\}$. A random $r$ is drawn uniformly from $[0, 1]$.
\begin{itemize}
\item If $q \geq r$: you win \$$H$ if red.
\item If $q < r$: you win \$$H$ with probability $r$.
\end{itemize}
Which $q$ do you pick?
\end{quote}

\vspace{0.3em}
\footnotesize
No payoff pairs shown. To find the maximizer, the subject must reason through what happens at each $r$ given each candidate $q$ --- full contingent reasoning demand preserved.
\end{frame}

\begin{frame}{Question 2: A, B, or both?}
The pure-incentives test diagnoses BIC failure by asking whether subjects identify the payoff-maximizer with beliefs stripped from the task.

\vspace{0.6em}
A and B differ in how much of the p-BDM structure remains:
\begin{itemize}
\item A: mechanism hidden; subjects see only event-contingent payoff pairs.
\item B: native p-BDM framing; contingent reasoning across $r$ preserved.
\end{itemize}

\vspace{0.4em}
Both designs are separate between-subject arms, with $\theta$ varying within-subject. Running A and B within-subject isolates menu-level vs. contingent-reasoning failure.

\vspace{0.6em}
\begin{block}{}
A, B, or both? Or a fourth design?
\end{block}
\end{frame}

\begin{frame}{MPL format: what the literature says}
Brown and Healy (2018) show that presentation format determines whether the MPL's IC survives subjects' behavior.

\vspace{0.7em}
List format induces subjects to treat the RPS list as one large decision and apply reduction of compound lotteries. ROCL combined with non-EU preferences forces a monotonicity violation (Segal 1990). Separated format (one row per screen, random order) does not trigger ROCL, so non-EU preferences have no mechanism-level consequence.

\vspace{0.7em}
Empirical result for risk elicitation: $p = 0.041$ (list) versus $p = 0.697$ (separated).
\end{frame}

\begin{frame}{MPL format: options and tradeoffs}
The MPL arm is the backbone of H2 (MPL $>$ BDM). For that test to be clean, the format has to preserve IC, yield interpretable responses, and keep burden tolerable.

\vspace{0.6em}
\small
\setlength{\tabcolsep}{4pt}
\begin{tabular}{@{}p{0.36\linewidth} p{0.13\linewidth} p{0.18\linewidth} p{0.18\linewidth}@{}}
\toprule
Format & Screens & IC defense & Burden \\
\midrule
Full list (all rows on 1 screen) & 1 & Weak & Low \\
Two-stage list (Holt \& Smith) & 2 & Weak & Moderate \\
Full separated (100 rows) & 100 & Strong & High \\
Coarse separated & 15--20 & Strong & Moderate \\
Two-stage separated & 20--30 & Strong & High \\
\bottomrule
\end{tabular}

\vspace{0.7em}
\footnotesize
List formats give precision at low burden but inherit the ROCL concern. Separated formats preserve IC but raise burden and can produce multi-switching. Two-stage variants trade extra screens for finer precision.
\end{frame}

\begin{frame}{Question 3: MPL format}
We are leaning toward coarse separated: 15 to 20 rows, one per screen, random order.

\vspace{0.7em}
This preserves IC per B\&H, keeps burden moderate, and gives 5pp precision, which is sufficient for our accuracy tests.

\vspace{0.8em}
\begin{block}{}
B\&H tested risk preferences with known lotteries in each row. Our event bet is ambiguous (subjective $\pi$). ROCL operates across rows; ambiguity operates within a row, and format does not close it. Does the cross-row transfer hold cleanly enough to commit, or do we need an auxiliary B\&H transfer arm (100--150 additional subjects)?
\end{block}
\end{frame}

\end{document}
